% !TEX root = ../main.tex
\chapter{Phân tích và thiết kế hệ thống}

\section{Phân tích tác nhân}

RabbitCV có bốn tác nhân chính:

\begin{itemize}
    \item \textbf{Khách:} chỉ truy cập tài nguyên công khai.
    \item \textbf{Ứng viên:} là tác nhân trung tâm của đề tài vì trực tiếp tạo CV, nhận gợi ý AI và dùng CV để ứng tuyển.
    \item \textbf{Doanh nghiệp:} tạo tin, tiếp nhận hồ sơ và trao đổi với ứng viên.
    \item \textbf{Quản trị viên:} giám sát dữ liệu và thực hiện các thao tác đặc quyền trong phạm vi backend đã triển khai.
\end{itemize}

Tác nhân ngoài hệ thống là dịch vụ AI Groq. Dịch vụ này không truy cập trực tiếp cơ sở dữ liệu; mọi yêu cầu đều đi qua backend để xác thực, rút gọn dữ liệu, áp dụng quota và yêu cầu đầu ra theo schema.

\begin{figure}[H]
    \centering
    \resizebox{0.96\textwidth}{!}{%
    \begin{tikzpicture}[
        actor/.style={draw, rounded corners, fill=gray!10, minimum width=2.2cm, minimum height=0.8cm, align=center},
        usecase/.style={draw, ellipse, fill=blue!6, minimum width=3.5cm, minimum height=0.82cm, align=center},
        link/.style={-Latex, line width=0.7pt}
    ]
        \node[actor] (guest) at (0,3.6) {Khách};
        \node[actor] (candidate) at (0,1.2) {Ứng viên};
        \node[actor] (company) at (12,3.6) {Doanh nghiệp};
        \node[actor] (admin) at (12,1.2) {Quản trị viên};
        \node[actor] (ai) at (12,-1.2) {Dịch vụ AI};

        \node[usecase] (public) at (4.1,4.2) {Xem và tìm việc công khai};
        \node[usecase] (auth) at (7.8,4.2) {Đăng ký, đăng nhập};
        \node[usecase] (cv) at (4.1,2.6) {Tạo, quản lý và in CV};
        \node[usecase] (aicv) at (7.8,2.6) {Viết, review và điền CV};
        \node[usecase] (apply) at (4.1,1.0) {Lưu việc và ứng tuyển};
        \node[usecase] (recruit) at (7.8,1.0) {Quản lý tin và hồ sơ};
        \node[usecase] (message) at (4.1,-0.6) {Chat và thông báo};
        \node[usecase] (score) at (7.8,-0.6) {Chấm CV, luyện phỏng vấn};
        \node[usecase] (manage) at (6.0,-2.0) {Quản trị dữ liệu hệ thống};

        \draw[link] (guest) -- (public);
        \draw[link] (guest) -- (auth);
        \draw[link] (candidate) -- (cv);
        \draw[link] (candidate) -- (aicv);
        \draw[link] (candidate) -- (apply);
        \draw[link] (candidate) -- (message);
        \draw[link] (company) -- (recruit);
        \draw[link] (company) -- (message);
        \draw[link] (admin) -- (manage);
        \draw[link] (ai) -- (aicv);
        \draw[link] (ai) -- (score);
        \node[draw, dashed, rounded corners, fit=(public)(auth)(cv)(aicv)(apply)(recruit)(message)(score)(manage), inner sep=0.35cm, label=above:{\textbf{RabbitCV}}] {};
    \end{tikzpicture}}
    \caption{Bản đồ tác nhân và nhóm ca sử dụng chính}
    \label{fig:usecase-map}
\end{figure}

\section{Yêu cầu chức năng}

Các yêu cầu được gắn mã để liên kết với module triển khai và kiểm thử.

\begin{longtable}{|p{0.13\textwidth}|p{0.20\textwidth}|p{0.57\textwidth}|}
    \hline
    \textbf{Mã} & \textbf{Nhóm} & \textbf{Yêu cầu} \\ \hline
    \endfirsthead
    \hline
    \textbf{Mã} & \textbf{Nhóm} & \textbf{Yêu cầu} \\ \hline
    \endhead
    FR-CV-01 & CV & Ứng viên tạo, đọc, cập nhật, xóa Resume và chọn template. \\ \hline
    FR-CV-02 & CV & Hệ thống preview CV theo A4, hỗ trợ in/tải PDF và tiếp tục chỉnh sửa dữ liệu cấu trúc. \\ \hline
    FR-CV-03 & PDF & Trong luồng ứng tuyển, ứng viên có thể tải CV PDF hợp lệ; doanh nghiệp xem từng trang bằng PDF.js. \\ \hline
    FR-AI-01 & AI viết CV & AI viết mới, cải thiện hoặc rút gọn phần tiểu sử; người dùng xem kết quả và quyết định áp dụng. \\ \hline
    FR-AI-02 & AI review/fill & AI review CV theo nhóm tiêu chí và gợi ý điền mục còn thiếu mà không tự khẳng định sự kiện cá nhân chưa có. \\ \hline
    FR-AI-03 & AI mở rộng & Hệ thống chấm mức độ phù hợp CV--Job và luyện phỏng vấn theo từng lượt. \\ \hline
    FR-JOB-01 & Việc làm & Người dùng tìm, lọc, xem, lưu và bỏ lưu tin tuyển dụng. \\ \hline
    FR-APP-01 & Ứng tuyển & Ứng viên nộp CV đã lưu hoặc PDF, không được nộp trùng và không nộp vào tin không khả dụng. \\ \hline
    FR-APP-02 & Lịch sử & Application giữ snapshot thông tin Job và trạng thái xử lý hồ sơ. \\ \hline
    FR-COM-01 & Doanh nghiệp & Doanh nghiệp tạo/sửa/xóa tin, quản lý nhận hồ sơ, xem ứng viên và cập nhật trạng thái. \\ \hline
    FR-MSG-01 & Giao tiếp & Hai bên gửi tin nhắn, theo dõi trạng thái đọc, ẩn room theo từng người và nhận thông báo nghiệp vụ. \\ \hline
    FR-ADM-01 & Quản trị & Admin xem thống kê và quản lý người dùng, doanh nghiệp, tin trong phạm vi API có sẵn. \\ \hline
\caption{Danh sách yêu cầu chức năng}
\label{tab:functional-requirements}
\end{longtable}

\section{Yêu cầu phi chức năng}

\begin{table}[H]
    \centering
    \small
    \begin{tabularx}{\textwidth}{|p{0.13\textwidth}|p{0.21\textwidth}|Y|} \hline
        \textbf{Mã} & \textbf{Nhóm} & \textbf{Tiêu chí kiểm chứng} \\ \hline
        NFR-SEC-01 & Xác thực & Mọi API riêng tư và toàn bộ route AI yêu cầu JWT hợp lệ; thao tác đặc quyền phải kiểm tra vai trò ở backend. API chung, đăng nhập, đăng ký và AI có giới hạn tần suất riêng. \\ \hline
        NFR-SEC-02 & Quyền dữ liệu & Resume chỉ được trả cho chủ sở hữu, admin hoặc doanh nghiệp có quan hệ nghiệp vụ phù hợp. \\ \hline
        NFR-PDF-01 & Tài liệu & CV web dùng kích thước A4; PDF upload phải có header hợp lệ và kích thước giải mã không quá 5 MB. \\ \hline
        NFR-AI-01 & Kiểm soát AI & Rate limit 30 request/15 phút; quota ngày và timeout có thể cấu hình, mặc định lần lượt là 100 request và 20 giây. \\ \hline
        NFR-AI-02 & Hợp đồng dữ liệu & Request AI loại thông tin liên hệ trực tiếp; response phải khớp JSON Schema và được làm sạch trước khi trả frontend. \\ \hline
        NFR-DATA-01 & Toàn vẹn & Unique index ngăn ứng tuyển trùng; snapshot giữ lịch sử; thao tác ẩn/đọc chat không xóa dữ liệu của bên còn lại. \\ \hline
        NFR-TEST-01 & Khả năng tái lập & Quy trình \texttt{npm run check} phải chạy ESLint, 41 test hiện có và Vite build; kết quả phải gắn với ngày, phiên bản runtime và trạng thái working tree. \\ \hline
        NFR-UX-01 & Khả dụng & Kết quả AI chỉ là gợi ý, có trạng thái chờ/lỗi và không tự ghi đè nội dung khi người dùng chưa xác nhận. \\ \hline
    \end{tabularx}
    \caption{Yêu cầu phi chức năng có tiêu chí kiểm chứng}
    \label{tab:nonfunctional-requirements}
\end{table}

\section{Ma trận phân quyền}

\begin{table}[H]
    \centering
    \scriptsize
    \begin{tabularx}{\textwidth}{|Y|c|c|c|c|} \hline
        \textbf{Chức năng} & \textbf{Khách} & \textbf{Ứng viên} & \textbf{Doanh nghiệp} & \textbf{Admin} \\ \hline
        Xem danh sách/chi tiết việc công khai & Có & Có & Có & Có \\ \hline
        Tạo, sửa và dùng AI cho CV cá nhân & Không & Có & Không & Giới hạn \\ \hline
        Lưu việc và nộp hồ sơ & Không & Có & Không & Không \\ \hline
        Tạo và quản lý tin tuyển dụng & Không & Không & Tin sở hữu & Đặc quyền \\ \hline
        Xem CV ứng viên & Không & CV sở hữu & Hồ sơ liên quan & Đặc quyền \\ \hline
        Cập nhật trạng thái Application & Không & Không & Hồ sơ thuộc tin sở hữu & Đặc quyền \\ \hline
        Nhắn tin và đọc thông báo & Không & Có & Có & Theo phạm vi \\ \hline
        Quản lý người dùng/doanh nghiệp & Không & Không & Không & Có \\ \hline
    \end{tabularx}
    \caption{Ma trận quyền theo tác nhân}
    \label{tab:role-matrix}
\end{table}

Ma trận trên mô tả chính sách mong muốn ở mức nghiệp vụ. Frontend có thể ẩn hoặc chuyển hướng giao diện, nhưng quyết định cuối cùng phải do middleware và truy vấn có điều kiện ở backend thực hiện. Qua đối chiếu mã nguồn hiện tại, nhóm API AI mới yêu cầu JWT chung và chưa kiểm tra riêng vai trò ứng viên; một số endpoint lưu việc và chat cũng mới dừng ở mức xác thực. Đây là khoảng cách giữa chính sách và hiện thực cần được kiểm thử, bổ sung trước khi triển khai thực tế.

\section{Kiến trúc tổng thể}

RabbitCV là ứng dụng client--server, không phải microservice. Frontend và backend có thể đóng gói/chạy riêng nhưng backend hiện tập trung phần lớn REST route trong \texttt{server/server.js}. Route AI, schema, service gọi nhà cung cấp, sanitizer và hàm xử lý hạn tin được tách thành các module riêng. Ở frontend, Context quản lý phiên và socket; TanStack Query quản lý một phần dữ liệu máy chủ; React Hook Form và Zod quản lý biểu mẫu xác thực.

\begin{figure}[H]
    \centering
    \resizebox{0.98\textwidth}{!}{%
    \begin{tikzpicture}[
        box/.style={draw, rounded corners, align=center, minimum height=1.05cm, minimum width=2.7cm, fill=blue!5},
        ext/.style={draw, rounded corners, align=center, minimum height=1.05cm, minimum width=2.5cm, fill=orange!10},
        arr/.style={-Latex, line width=0.8pt}
    ]
        \node[box] (browser) {Trình duyệt\\React + Router};
        \node[box, right=1.1cm of browser] (contexts) {Context + Query/Form\\Tailwind + Tiptap};
        \node[box, right=1.1cm of contexts] (api) {Express REST API\\JWT + rate limit};
        \node[box, right=1.1cm of api] (models) {Service/Utility\\Mongoose models};
        \node[box, right=1.1cm of models] (db) {MongoDB};
        \node[ext, below=1.1cm of api] (socket) {Socket.IO\\room người dùng};
        \node[ext, below=1.1cm of models] (cron) {Vercel Cron\\HTTP GET có token};
        \node[ext, above=1.1cm of models] (groq) {Groq API\\LLM + JSON Schema};
        \node[ext, above=1.1cm of contexts] (pdf) {PDF.js\\trình kết xuất};

        \draw[arr,<->] (browser) -- node[above]{HTTP} (contexts);
        \draw[arr,<->] (contexts) -- node[above]{JSON} (api);
        \draw[arr,<->] (api) -- (models);
        \draw[arr,<->] (models) -- (db);
        \draw[arr,<->] (contexts) -- (socket);
        \draw[arr,<->] (socket) -- (api);
        \draw[arr] (cron) -- (api);
        \draw[arr,<->] (api) -- (groq);
        \draw[arr] (contexts) -- (pdf);
    \end{tikzpicture}}
    \caption{Kiến trúc triển khai của RabbitCV}
    \label{fig:system-architecture}
\end{figure}

\section{Thiết kế mô-đun}

\begin{table}[H]
    \centering
    \small
    \begin{tabularx}{\textwidth}{|p{0.20\textwidth}|p{0.28\textwidth}|Y|} \hline
        \textbf{Mô-đun} & \textbf{Thành phần chính} & \textbf{Trách nhiệm} \\ \hline
        Phiên và vai trò & AuthContext, RoleAccessGuard, authMiddleware. & Đăng ký, đăng nhập, khôi phục phiên, đăng xuất, chuyển hướng và kiểm tra quyền. \\ \hline
        Dữ liệu giao diện & QueryClient, SavedJobsContext, useNotifications, form schema. & Cache dữ liệu máy chủ, polling, cập nhật lạc quan, khôi phục khi lỗi và kiểm tra biểu mẫu. \\ \hline
        CV & ResumeForm, ResumePreview, Resume model. & Dữ liệu CV, template, preview A4, in PDF, CV upload và quyền sở hữu. \\ \hline
        AI & aiRoutes, aiService, aiSchemas, sanitizer, AiUsage. & Chuẩn hóa input, quota, gọi model, kiểm tra schema, lọc output và thống kê. \\ \hline
        Việc làm & Jobs, JobDetail, Job/Application models. & Tìm kiếm, trạng thái tin, lưu việc, ứng tuyển, snapshot và xử lý hồ sơ. \\ \hline
        Giao tiếp & Message, SocketContext, NotificationBell. & Lưu message, phát event, trạng thái đọc/ẩn room và thông báo. \\ \hline
        Doanh nghiệp & CompanyDashboard và API Job/Application. & Hồ sơ công ty, tin tuyển dụng, danh sách ứng viên và cập nhật trạng thái. \\ \hline
        Quản trị & AdminLayout và nhóm API admin. & Thống kê, người dùng, doanh nghiệp và tin tuyển dụng. \\ \hline
    \end{tabularx}
    \caption{Phân rã mô-đun theo thành phần triển khai}
    \label{tab:module-design}
\end{table}

\section{Thiết kế dữ liệu}

MongoDB dùng reference ObjectId giữa các tài liệu thay vì lược đồ quan hệ vật lý. Hình \ref{fig:data-model} thể hiện quan hệ nghiệp vụ chính; các mũi tên không có nghĩa database tự áp dụng khóa ngoại.

\begin{figure}[H]
    \centering
    \resizebox{0.96\textwidth}{!}{%
    \begin{tikzpicture}[
        entity/.style={draw, rounded corners, align=left, text width=3.15cm, inner sep=4pt, minimum height=1.05cm, fill=green!6, font=\small},
        rel/.style={-Latex, line width=0.75pt}
    ]
        \node[entity] (user) at (5.8,6.0) {\textbf{User}\\role, profile, savedJobs};
        \node[entity] (msg) at (-0.2,3.0) {\textbf{Message}\\sender, room, content};
        \node[entity] (resume) at (3.8,3.0) {\textbf{Resume}\\userId, sections, pdfData};
        \node[entity] (job) at (7.8,3.0) {\textbf{Job}\\companyId, status, expiryDate};
        \node[entity] (usage) at (11.8,3.0) {\textbf{AiUsage}\\userId, feature, model\\latency, token, status};
        \node[entity] (read) at (0.4,0) {\textbf{ChatReadState /}\\\textbf{HiddenChat}\\userId, room, lastReadAt};
        \node[entity] (app) at (5.8,0) {\textbf{Application}\\jobId, candidateId, resumeId\\jobSnapshot, status};
        \node[entity] (noti) at (11.0,0) {\textbf{Notification}\\recipient, type, isRead};

        \draw[rel] (user) -- (msg);
        \draw[rel] (user) -- (resume);
        \draw[rel] (user) -- (job);
        \draw[rel] (user) -- (usage);
        \draw[rel] (user) -- (app);
        \draw[rel] (resume) -- (app);
        \draw[rel] (job) -- (app);
        \draw[rel] (msg) -- (read);
        \draw[rel] (app) -- (noti);
    \end{tikzpicture}}
    \caption{Mô hình quan hệ giữa các collection chính}
    \label{fig:data-model}
\end{figure}

\subsection{User và Resume}

User lưu username, email, mật khẩu băm, vai trò, trạng thái tài khoản, hồ sơ ứng viên hoặc công ty và danh sách Job đã lưu. Resume tham chiếu \texttt{userId}, chứa dữ liệu CV có cấu trúc, mẫu và metadata. \texttt{isVirtual} phân biệt PDF tạm được tạo khi ứng tuyển với CV cấu trúc do người dùng quản lý; \texttt{pdfData} chỉ có ý nghĩa với CV PDF và bị loại khỏi truy vấn danh sách.

\subsection{Job và Application}

Job liên kết doanh nghiệp qua \texttt{companyId}, chứa nội dung, trạng thái và ngày hết hạn. Application liên kết Job, ứng viên và Resume. Unique index trên \texttt{jobId}--\texttt{candidateId} ngăn hồ sơ trùng. \texttt{jobSnapshot} lưu tiêu đề, công ty, địa điểm, loại và lương tại thời điểm nộp để lịch sử không mất ngữ cảnh khi Job bị xóa.

\subsection{Message, Notification và trạng thái đọc}

Message lưu sender, room và nội dung. ChatReadState lưu \texttt{lastReadAt} theo cặp user--room; HiddenChat ghi việc một phía ẩn room. Notification lưu recipient, nội dung, liên kết và \texttt{isRead}. Thiết kế này tách trạng thái cá nhân khỏi dữ liệu trao đổi chung.

\subsection{AiUsage}

AiUsage lưu user, feature, model, hash đầu vào, trạng thái, token, độ trễ và mã lỗi. Dữ liệu này hỗ trợ quota và quan sát kỹ thuật mà không cần lưu nguyên prompt hoặc toàn bộ CV. Hash không được dùng như cơ chế ẩn danh tuyệt đối; hệ thống vẫn cần hạn chế quyền truy cập collection thống kê.

\section{Các quy trình nghiệp vụ chính}

\subsection{Ứng tuyển bằng CV}

\begin{figure}[H]
    \centering
    \resizebox{0.96\textwidth}{!}{%
    \begin{tikzpicture}[
        flowstep/.style={draw, rounded corners, align=center, minimum width=2.65cm, minimum height=1.05cm, fill=purple!7},
        arr/.style={-Latex, line width=0.8pt}
    ]
        \node[flowstep] (select) {Ứng viên chọn\\Job và CV};
        \node[flowstep, right=0.65cm of select] (auth) {Xác thực và\\quyền sở hữu CV};
        \node[flowstep, right=0.65cm of auth] (check) {Kiểm tra Job\\hạn và trạng thái};
        \node[flowstep, right=0.65cm of check] (unique) {Kiểm tra unique\\Application};
        \node[flowstep, right=0.65cm of unique] (save) {Lưu Application\\và jobSnapshot};
        \node[flowstep, right=0.65cm of save] (notify) {Tạo thông báo\\trả kết quả};
        \draw[arr] (select) -- (auth);
        \draw[arr] (auth) -- (check);
        \draw[arr] (check) -- (unique);
        \draw[arr] (unique) -- (save);
        \draw[arr] (save) -- (notify);
    \end{tikzpicture}}
    \caption{Luồng kiểm soát khi ứng tuyển bằng CV}
    \label{fig:application-flow}
\end{figure}

Với PDF upload, backend kiểm tra data URL, MIME, đuôi, header và dung lượng giải mã; sau đó tạo Resume ảo và tiếp tục cùng luồng Application. Nếu bất kỳ kiểm tra nào thất bại, không tạo Resume hoặc Application mới.

\subsection{Yêu cầu AI}

\begin{figure}[H]
    \centering
    \resizebox{0.98\textwidth}{!}{%
    \begin{tikzpicture}[
        flowstep/.style={draw, rounded corners, align=center, minimum width=2.45cm, minimum height=1.05cm, fill=orange!9},
        arr/.style={-Latex, line width=0.8pt}
    ]
        \node[flowstep] (ui) {UI gửi feature\\và context};
        \node[flowstep, right=0.5cm of ui] (gate) {JWT, rate limit\\quota ngày};
        \node[flowstep, right=0.5cm of gate] (clean) {Sanitize và\\rút gọn dữ liệu};
        \node[flowstep, right=0.5cm of clean] (schema) {Chọn instruction\\và JSON Schema};
        \node[flowstep, right=0.5cm of schema] (provider) {Groq LLM\\structured output};
        \node[flowstep, right=0.5cm of provider] (validate) {Parse, kiểm tra\\nghiệp vụ và lọc};
        \node[flowstep, right=0.5cm of validate] (result) {Gợi ý + meta\\người dùng áp dụng};
        \draw[arr] (ui) -- (gate);
        \draw[arr] (gate) -- (clean);
        \draw[arr] (clean) -- (schema);
        \draw[arr] (schema) -- (provider);
        \draw[arr] (provider) -- (validate);
        \draw[arr] (validate) -- (result);
    \end{tikzpicture}}
    \caption{Đường đi của dữ liệu trong một yêu cầu AI}
    \label{fig:ai-request-flow}
\end{figure}

Service tạo bản ghi AiUsage trước khi gọi nhà cung cấp và cập nhật trạng thái sau khi hoàn tất hoặc thất bại. Nhà cung cấp được yêu cầu áp dụng JSON Schema nghiêm ngặt; backend tiếp tục kiểm tra điều kiện ngữ nghĩa ở luồng tiểu sử và làm sạch/giới hạn đầu ra ở luồng điền CV, phỏng vấn. Frontend nhận lưu ý cùng kết quả ở các chức năng review, chấm CV và phỏng vấn. Không có bước tự động lưu kết quả AI vào Resume.

\subsection{Chat và trạng thái đọc}

Người gửi gọi REST API; server lấy định danh và tên người gửi từ tài khoản đã xác thực, kiểm tra tin có nội dung hoặc ảnh, lưu Message rồi phát \texttt{newMessage} đến phòng của người nhận. Văn bản chat được React hiển thị như chuỗi, còn ảnh được nén ở client trước khi chuyển thành data URL. Khi mở phòng, client cập nhật \texttt{lastReadAt}. Việc ẩn phòng tạo hoặc cập nhật HiddenChat theo user; Message của phía còn lại không bị xóa. Notification không dùng chung kênh này mà được tải qua API, do đó có thể có độ trễ tối đa theo chu kỳ polling.

\section{Bảo mật và ranh giới tin cậy}

Thiết kế xác định ba ranh giới tin cậy chính:

\begin{itemize}
    \item \textbf{Giữa trình duyệt và backend:} mọi dữ liệu từ client đều không đáng tin, kể cả vai trò hoặc mã người dùng.
    \item \textbf{Giữa backend và dịch vụ AI:} chỉ dữ liệu đã được làm sạch, giới hạn và thật sự cần thiết mới được gửi ra ngoài.
    \item \textbf{Tại dữ liệu CV:} đường dẫn xem CV phải kiểm tra quan hệ sở hữu hoặc quan hệ tuyển dụng, không chỉ kiểm tra người gọi đã đăng nhập.
\end{itemize}

Các biện pháp kiểm soát hiện có được chia thành các nhóm sau:

\begin{itemize}
    \item \textbf{Xác thực và bảo vệ HTTP:} bcryptjs, JWT trong cookie HTTP-only, Helmet, CSP, CORS theo danh sách nguồn và giới hạn tốc độ.
    \item \textbf{Kiểm tra dữ liệu và quyền truy cập:} kiểm tra ObjectId, danh sách trường người dùng theo ngữ cảnh, bộ làm sạch rich text/PDF/AI, chỉ mục duy nhất và truy vấn kèm điều kiện sở hữu.
    \item \textbf{Giới hạn thông tin vận hành:} health check công khai chỉ trả trạng thái tối giản; thông tin môi trường và cơ sở dữ liệu chi tiết yêu cầu quyền quản trị.
    \item \textbf{Phạm vi bảo vệ:} các biện pháp trên giúp giảm rủi ro nhưng không thay thế kiểm thử xâm nhập, phân tích CSRF, quản lý bí mật production hoặc rà soát đầy đủ ma trận quyền.
\end{itemize}

\noindent\textbf{Hạn chế của thiết kế hiện tại.}
Backend còn tập trung nhiều route trong một file; PDF và ảnh chat dạng data URL làm tài liệu MongoDB lớn; notification dùng polling; Socket.IO trên môi trường serverless chưa phải kiến trúc phù hợp cho tải lớn. Tác vụ theo lịch chưa có kiểm thử về gọi đồng thời, retry và quan sát lỗi. Mã nguồn vẫn có khóa JWT dự phòng và cơ chế tạo tài khoản quản trị phục vụ trình diễn; chúng phải bị vô hiệu hóa hoặc bắt buộc cấu hình ở production. Đặc biệt, API gửi/ẩn chat chưa kiểm tra đầy đủ người gọi có thuộc phòng và API AI chưa chặn theo vai trò ứng viên. Phần AI còn phụ thuộc nhà cung cấp và model cấu hình, trong khi chất lượng ngữ nghĩa chưa có bộ dữ liệu chấm độc lập. Các hạn chế này được chuyển thành mục kiểm thử hoặc hướng phát triển thay vì được mô tả như chức năng đã hoàn tất.
